\documentclass[12pt]{article}

\usepackage[margin=1in]{geometry}

\usepackage{amsmath,amssymb,booktabs,longtable,array,hyperref}

\usepackage{newtxtext,newtxmath}

\title{Matrix-Free Second-Order Optimization of Nonorthogonal Valence-Bond Wave Functions}

\author{[Author names and affiliations to be confirmed]}

\date{Internal review draft: 7 October 2026}

\begin{document}

\maketitle

% REVIEW: 暂定题目。先以非正交 VBSCF 的 matrix-free 二阶框架为主线；是否将 NEO 作为主方法、TNHVP 作为对照，需结合统一版本的测试结果决定。

JCTC Article --- manuscript outline for internal review, 7 October 2026


\section*{ABSTRACT}

Orbital optimization in valence-bond self-consistent-field theory must accommodate nonorthogonal orbitals, prescribed atomic-orbital supports, and relaxation of the structure coefficients. We formulate a matrix-free second-order treatment in support-preserving nonredundant coordinates and distinguish the relaxed-orbital and coupled orbital--structure models. The numerical study will compare truncated Newton and norm-extended optimization with conventional and block-preconditioned L-BFGS, reporting convergence, operator work, wall time, and peak memory under matched conditions. [Replace this sentence with the principal quantitative findings and the supported conclusion after the benchmark campaign.]


% REVIEW: 这是四句摘要框架，不是已完成研究的结果摘要。不能把‘外层步数减少’直接写成‘总计算更快’，也不能预先承诺所有体系都二次收敛。

\section{INTRODUCTION}

\subsection{The optimization problem}

Explain why VBSCF uses nonorthogonal localized orbitals and, in HAO calculations, fixed orbital supports. Separate the cost of optimizing orbitals from the cost of solving the structure-coefficient problem. Introduce OEO as full-AO support, not as an orthogonal-orbital model.


\subsection{Relation to second-order MCSCF optimization}

Review restricted-step Newton, augmented-Hessian and norm-extended optimization, and matrix-free CASSCF implementations [1--3]. Identify what carries over to VBSCF---the local coupled quadratic model and iterative operator actions---and what does not carry over without derivation: orthogonal orbital rotations and their redundancy rules.


\subsection{Scope of this work}

State the methodological contribution as a second-order optimization framework for the specified nonorthogonal VB variational space. Treat integral representation and string contraction as supporting implementations, not independent novelty claims in this paper. The principal numerical question is when additional curvature information improves convergence enough to repay its cost relative to block-LBFGS.


\section{THEORY}

\subsection{Stationary structure problem and state averaging}

Define the generalized structure eigenproblem at fixed orbital coefficients x. Use a different symbol for the AO overlap and the structure overlap throughout the manuscript.


\begin{equation}

\mathbf H_{\mathrm{VB}}(\mathbf x)\mathbf a_s=E_s\mathbf S_{\mathrm{VB}}(\mathbf x)\mathbf a_s,\quad\mathbf a_r^T\mathbf S_{\mathrm{VB}}\mathbf a_s=\delta_{rs}

\end{equation}

\begin{equation}

\mathcal E(\mathbf x)=\sum_{s\in\mathcal A}w_sE_s(\mathbf x),\quad w_s=1/m

\end{equation}

For the single-state case, use the lowest stationary eigenvalue; for equal-weight state averaging, define the selected eigenspace and its separation from excluded roots. Derive normalization and selected-subspace gauge projections before introducing response equations.


\subsection{Support-preserving nonredundant orbital coordinates}

\begin{equation}

\mathbf c_p=\mathbf R_p\mathbf x_p,\quad \mathbf x(\mathbf p)=\mathcal R_k(\mathbf p),\quad\mathcal R_k(\mathbf 0)=\mathbf x_k

\end{equation}

Define the support injection R_p, horizontal lift U_k, local update map, and the physical step metric M_o. Classify a direction as redundant only if it leaves the represented wave function unchanged within the selected structure space. In particular, active--active mixing is not automatically redundant for an arbitrary selected-VB expansion.


\begin{equation}

\mathbf g_o=\nabla_{\mathbf p}(\mathcal E\circ\mathcal R_k)(\mathbf 0),\quad \mathbf H_o\mathbf v=D^2(\mathcal E\circ\mathcal R_k)(\mathbf 0)[\mathbf v]

\end{equation}

Give the derivative of the actual update map, rather than transferring an orbital-rotation Hessian to a coefficient update. Check the consistency of lift, norm, and retraction.


% REVIEW: 最终稿需从冻结的代码版本核对 retraction 和 pullback。历史理论文档中的 additive chart 不能不加检查地当作当前实现。

\subsection{Orbital--structure coupling and the relaxed Hessian}

At a structure-stationary accepted point, write the constrained local Hessian in orbital and projected structure tangent variables. A is the fixed-structure orbital block, B is the mixed derivative, and C is the structure tangent block. All blocks include the chosen normalization and generalized-overlap constraints.


\begin{equation}

\mathbf z=\begin{pmatrix}\mathbf p\\\mathbf q\end{pmatrix},\quad\mathbf g=\begin{pmatrix}\mathbf g_o\\\mathbf 0\end{pmatrix},\quad\mathbf K=\begin{pmatrix}\mathbf A&\mathbf B^T\\\mathbf B&\mathbf C\end{pmatrix}

\end{equation}

\begin{equation}

\mathbf H_{\mathrm{rel}}=\mathbf A-\mathbf B^T\mathbf C^\dagger\mathbf B,\quad\mathbf q_N=-\mathbf C^\dagger\mathbf B\mathbf p_N

\end{equation}

Define \(C^\dagger\) only on the admissible response space. Explain the equivalence of the unshifted coupled Newton and relaxed-orbital Newton limits, and why shifted coupled and eliminated models need not be identical away from that limit.


\subsection{Restricted steps and residual certification}

\begin{equation}

m(\mathbf z)=\mathcal E_k+\mathbf g^T\mathbf z+\tfrac12\mathbf z^T\mathbf K\mathbf z,\quad\mathbf z^T\mathbf M\mathbf z\leq\Delta_k^2

\end{equation}

\begin{equation}

\mathbf r=(\mathbf K+\lambda\mathbf M)\mathbf z+\mathbf g,\quad\lambda\geq0,\quad\mathbf K+\lambda\mathbf M\succeq0

\end{equation}

\begin{equation}

\lambda(\mathbf z^T\mathbf M\mathbf z-\Delta_k^2)=0,\quad\|\mathbf r\|_{\mathbf M^{-1}}\leq\epsilon_{\mathrm{abs}}+\eta_k\|\mathbf g\|_{\mathbf M^{-1}}

\end{equation}

Separate stationarity, step feasibility, complementarity, and curvature checks. State the regularity and inner-solve conditions required for superlinear or quadratic local convergence; an HVP implementation alone does not establish those properties.


\subsection{TNHVP and norm-extended optimization}

Describe TNHVP as a truncated Newton realization of the relaxed model, and NEO as a norm-extended formulation of the local second-order model [1]. Define the gradient scaling and the step extracted from the generalized augmented eigenvector explicitly.


\begin{equation}

\begin{pmatrix}0&\alpha\mathbf g^T\\\alpha\mathbf g&\mathbf K\end{pmatrix}\begin{pmatrix}\beta\\\mathbf y\end{pmatrix}=\mu\begin{pmatrix}1&0\\0&\mathbf M\end{pmatrix}\begin{pmatrix}\beta\\\mathbf y\end{pmatrix}

\end{equation}

\begin{equation}

\mathbf z=\mathbf y/(\alpha\beta),\quad\lambda=-\mu,\quad(\mathbf K+\lambda\mathbf M)\mathbf z=-\mathbf g

\end{equation}

Discuss the nonzero homogeneous component, lowest-root selection, and the interior Newton limit. Describe how the projected solver reaches this model; do not equate a generic augmented-Hessian step with a certified norm-constrained solution.


\subsection{Preconditioning and controlled structure response}

Distinguish the Hessian model from the preconditioner. Present the one-electron block preconditioner, retained secant information where applicable, and the role of coupled structure response. If the implementation omits or restricts response according to a work estimate, identify that calculation as an approximate model and report it separately.


% REVIEW: 当前 NEO 后端仍有 structure-response admission 分支。正文不能把所有生产运行都描述为完整 coupled Newton；必须分别记录 admitted / core-only 路径。

\section{IMPLEMENTATION}

\subsection{Matrix-free derivative actions}

Describe the accepted-point data required for gradient and HVP evaluation: normalization, inactive projection, integral transformation, structure action, and response. Identify reused quantities and directional quantities, without claiming that matrix-free means no large intermediates. Give memory costs in orbital, structure, and unique-string dimensions.


\subsection{Trial steps and accepted-point updates}

\begin{equation}

\rho_k=\frac{\mathcal E_k-\mathcal E(\mathcal R_k(\mathbf p))}{m(\mathbf0)-m(\mathbf z)}

\end{equation}

Explain trial evaluation, acceptance, trust-radius updates, and rebuilding the local model. State which quantities are refreshed at a keyframe and which are reused during subspace expansion. Count rejected trials as work even when they do not advance the accepted-iteration counter.


Algorithm 1. Accepted-point second-order VBSCF optimization (outline)


1. Read the same orbital guess and structure list used by all comparison methods.


2. Solve the selected structure problem; evaluate energy and the physical orbital gradient.


3. Build the support-preserving local coordinates, metric, derivative actions, and preconditioner.


4. Expand the iterative second-order model; certify its residual and record response treatment.


5. Retract the orbital step, solve the trial structure problem, and evaluate model agreement.


6. Accept or reject the trial; refresh accepted-point data and test both stopping criteria.


\subsection{Cost accounting}

Report orbital HVPs, orbital-to-structure and transpose-coupling actions, structure H/S actions, trial evaluations, and accepted steps separately. Include setup, final observables, and total process memory in end-to-end measurements; isolate optimizer time in an additional breakdown.


\section{COMPUTATIONAL DETAILS}

\subsection{Methods and matched conditions}

\begin{longtable}{p{0.3\textwidth}|p{0.3\textwidth}|p{0.3\textwidth}}

\toprule

Method & Selector & Role \\

\midrule

\endhead

XMVB-style L-BFGS & ISCF=5 & Conventional quasi-Newton baseline \\

TNHVP & ISCF=7 & Matrix-free truncated Newton comparison \\

Block-LBFGS & ISCF=8 & Enhanced baseline with one-electron block initial inverse \\

NEO & ISCF=9 & Norm-extended second-order method \\

\bottomrule

\end{longtable}

Table 1. Comparison methods. Input selectors are implementation details, not method names.


Use identical initial orbitals, geometry, basis, active electrons/orbitals, spin, selected structure list, state weights, integral representation, and final convergence tests. Compare exact and RI only in separate matched groups. Dense and Davidson structure solvers must solve the same selected generalized eigenproblem to controlled accuracy.


Record the projected-gradient convention and the physical dual-metric norm. The working energy-change and projected-gradient targets are \(10^{-7}\) hartree and \(10^{-3}\), respectively; confirm the precise implemented norms and final settings before collection. Use an independently tighter calculation for endpoint comparisons.


% REVIEW: 用户此前认可的阈值作为工作设置；不能把不同坐标中的无穷范数当作完全等价的物理梯度。正式稿应报告实际采用的定义。

\subsection{Test-set design}

\begin{longtable}{p{0.3\textwidth}|p{0.3\textwidth}|p{0.3\textwidth}}

\toprule

Group & Examples & Purpose \\

\midrule

\endhead

Small derivative oracles & F\(_2\); small selected-structure cases & Finite differences, dense Hessian, endpoint checks \\

Localized HAO & 241 (benzene), 240, 7975, 7963, YAMSAI & Support constraints and convergence variation \\

Transition metals & MnF\(_2\) and FeCl\(_2\) & Conditioning, spin, and OEO/HAO behavior \\

Larger response spaces & CERRAS and LOFLEA & Response cost and memory \\

Equal-weight state average & Small multistate cases; selected pyrrole/OEO cases & Root-space and response stability \\

\bottomrule

\end{longtable}

Table 2. Proposed test groups; molecular specifications and eligibility remain to be frozen.


Do not include unresolved endpoint discrepancies as successful validation. A PySCF CASSCF comparison is a variational-space reference only when the OEO structure expansion, spin/symmetry sector, integrals, and core convention match; otherwise it is a diagnostic comparison.


% REVIEW: FeCl₂、部分 pyrrole/OEO 的历史问题要在版本冻结时重新确认。这里仅列测试候选，不把历史能量差当成已解决。

\subsection{Hardware and reproducibility}

Collect publication timings on Hanhai25 compute nodes through Slurm. Report CPU model, allocated cores, affinity, OpenMP/BLAS threads, compiler flags, dependency versions, and commit identifier. Record peak RSS at the job/process level. Preserve nonconverged runs and resource failures in the dataset rather than dropping them.


\section{RESULTS AND DISCUSSION}

\subsection{Derivative and response validation}

Table 3 will report gradient directional differences, HVP directional differences, bilinear symmetry defects, and coupled-versus-Schur agreement on small cases. Include nonorthogonal well-conditioned and near-singular orbital points, HAO and OEO, and at least one equal-weight state average. [Insert verified numerical errors.]


\subsection{Local convergence}

Figure 2 will show energy error and physical gradient norm for identical controlled perturbations around a common stationary point. Establish the local convergence regime separately from the number of iterations required to enter it. [Insert trajectories and the residual forcing sequence.]


\subsection{End-to-end comparison}

Figures 3 and 4 will plot accepted steps and total wall time against a common system ordering, with failures visibly marked. Table 4 will provide final energies, energy changes, physical gradients, action counts, wall time, and RSS. Discuss cases where fewer steps do not yield lower wall time. [Insert matched results; do not mix revisions.]


\subsection{Where the second-order cost is spent}

Figure 5 will break down setup, core HVP, coupling/structure response, trial evaluation, and final output. Analyze larger active spaces using measured work and memory rather than accepted-step count alone. Compare full and restricted response with the model treatment stated for every run. [Insert stage timings and response statistics.]


\subsection{Ablations and limitations}

Separate model changes from preconditioner changes. Prioritize matched full/restricted response and one-electron block/Hessian-diagonal comparisons; identity preconditioning need not be a production benchmark. Discuss small structure gaps, incomplete selected-VB spaces, root changes, and redundant-coordinate conditioning. [Insert only ablations that answer a specific algorithmic question.]


\section{CONCLUSIONS}

Summarize the validated second-order formulation, the measured range of efficiency, and the cost of structure relaxation. State which method is recommended for which workload only if the data support that distinction. [Write the quantitative conclusion after Sections 5.1--5.5 are complete.]


\section*{ASSOCIATED CONTENT}

\subsection{Supporting Information}

Planned contents: coordinate and response derivations; update-map derivatives; finite-difference protocols; complete molecular specifications and initial orbitals; raw convergence trajectories; nonconvergence records; timing scripts and environment details. Submit the SI separately from the main manuscript.


\section*{AUTHOR INFORMATION}

Corresponding Author: [Name and e-mail]. Author Contributions: [To be agreed by all authors]. Notes: [Competing-interest statement to be confirmed].


\section*{ACKNOWLEDGMENTS}

[Funding, computing resources, and individual acknowledgments to be confirmed. Add any required AI-assistance disclosure after review of the applicable ACS policy; do not list an AI tool as an author.]


\section*{DATA AVAILABILITY STATEMENT}

[Supply a persistent public repository link or DOI for the input files, initial orbitals, raw outputs, benchmark metadata, and analysis scripts. Identify the code release and commit used for every figure. Explain any justified restrictions.]


\section*{REFERENCES --- STARTER LIST}

(1) Jensen, H.-J. Aa.; Jørgensen, P. A Direct Approach to Second-Order MCSCF Calculations Using a Norm Extended Optimization Scheme. J. Chem. Phys. 1984, 80, 1204--1214. https://doi.org/10.1063/1.446797.


(2) Werner, H.-J.; Knowles, P. J. A Second Order Multiconfiguration SCF Procedure with Optimum Convergence. J. Chem. Phys. 1985, 82, 5053--5063. https://doi.org/10.1063/1.448627.


(3) Sun, Q.; Yang, J.; Chan, G. K.-L. A General Second Order Complete Active Space Self-Consistent-Field Solver for Large-Scale Systems. Chem. Phys. Lett. 2017, 683, 291--299. https://doi.org/10.1016/j.cplett.2017.03.004.


[Complete the VBSCF/XMVB, inexact-Newton, trust-region, L-BFGS, and density-fitting literature before submission. This starter list is not an exhaustive priority review.]


\clearpage

\section*{INTERNAL DRAFT NOTES --- REMOVE BEFORE SUBMISSION}

This is a LaTeX master outline, not a finished manuscript. Resolve placeholders before submission. The Word review snapshot is generated with Pandoc; equations are editable Word mathematics. Review comments from earlier snapshots are not automatically transferred.


The LaTeX master uses the standard article class, not an official ACS LaTeX class. The Word snapshot takes its formatting from the separately retained ACS-based Word draft. This experiment does not establish submission readiness.


JCTC author guidelines: https://researcher-resources.acs.org/publish/author_guidelines?coden=jctcce


Official ACS Word template: https://pubsapp.acs.org/paragonplus/submission/templates/acs_template_msw2010_ol_pc.dotx


Code mapping checked for this outline: ISCF=5 conventional XMVB-style L-BFGS; ISCF=7 TNHVP; ISCF=8 block-LBFGS; ISCF=9 NEO. Historical theory notes are not treated as proof of current implementation behavior. No new molecular calculations were performed.


Editorial decisions: (i) freeze the primary method and exact response policy; (ii) resolve endpoint validation for difficult OEO cases; (iii) freeze the benchmark version and matched inputs; (iv) select figures only after complete data are available; (v) agree authorship, funding, data deposition, and final novelty claims.


\end{document}
